\section{Wong Simulations}

% latex table generated in R 4.4.2 by xtable 1.8-4 package
% Fri Jan  9 02:38:07 2026
\begin{table}[ht]
\centering
\begin{tabular}{rrrrrrrrrrrrrrrrrrrrr}
  \hline
 & ar & -0.9 & -0.8 & -0.7 & -0.6 & -0.5 & -0.4 & -0.3 & -0.2 & -0.1 & 0 & 0.1 & 0.2 & 0.3 & 0.4 & 0.5 & 0.6 & 0.7 & 0.8 & 0.9 \\ 
  \hline
1 & -0.90 & 0.95 & 0.96 & 0.96 & 0.96 & 0.95 & 0.95 & 0.96 & 0.95 & 0.95 & 0.94 & 0.94 & 0.95 & 0.96 & 0.95 & 0.95 & 0.95 & 0.95 & 0.96 &  \\ 
  2 & -0.80 & 0.94 & 0.94 & 0.95 & 0.94 & 0.94 &  & 0.95 & 0.95 & 0.95 & 0.95 & 0.95 & 0.96 & 0.95 & 0.95 & 0.95 & 0.95 & 0.93 &  & 0.93 \\ 
  3 & -0.70 & 0.95 & 0.96 & 0.96 & 0.95 & 0.95 & 0.96 & 0.95 & 0.95 & 0.94 & 0.96 & 0.96 & 0.95 & 0.95 & 0.94 & 0.94 & 0.94 &  & 0.94 & 0.95 \\ 
  4 & -0.60 & 0.96 & 0.95 & 0.93 & 0.95 & 0.94 & 0.95 & 0.00 & 0.94 & 0.94 & 0.95 & 0.95 & 0.95 & 0.95 & 0.95 & 0.93 &  & 0.94 & 0.95 & 0.95 \\ 
  5 & -0.50 & 0.95 & 0.94 & 0.95 & 0.95 & 0.96 & 0.95 & 0.95 & 0.95 & 0.94 & 0.95 & 0.94 & 0.94 & 0.94 & 0.93 &  & 0.93 & 0.94 & 0.94 & 0.94 \\ 
  6 & -0.40 & 0.95 & 0.96 & 0.95 & 0.95 & 0.95 & 0.95 & 0.96 &  & 0.95 & 0.96 & 0.96 & 0.96 & 0.92 &  & 0.93 & 0.95 & 0.96 & 0.95 & 0.95 \\ 
  7 & -0.30 & 0.95 & 0.96 & 0.94 & 0.95 & 0.95 & 0.95 & 0.96 & 0.96 & 0.94 & 0.95 & 0.94 & 0.95 &  & 0.94 & 0.95 & 0.95 & 0.96 & 0.96 & 0.94 \\ 
  8 & -0.20 & 0.95 & 0.95 & 0.95 & 0.94 & 0.95 & 0.95 & 0.95 & 0.95 &  & 0.95 & 0.93 &  & 0.95 & 0.94 & 0.96 & 0.95 & 0.95 & 0.94 & 0.95 \\ 
  9 & -0.10 & 0.94 & 0.94 & 0.94 & 0.97 & 0.94 & 0.95 & 0.95 & 0.94 & 0.93 & 0.94 &  & 0.95 & 0.95 & 0.95 & 0.96 & 0.94 & 0.95 & 0.97 & 0.95 \\ 
  10 & 0.00 & 0.95 & 0.95 & 0.95 & 0.95 & 0.96 & 0.95 & 0.95 & 0.95 & 0.93 &  & 0.95 & 0.95 & 0.96 & 0.95 & 0.95 & 0.95 & 0.94 & 0.95 & 0.94 \\ 
  11 & 0.10 & 0.95 & 0.95 & 0.96 & 0.96 & 0.96 & 0.93 & 0.94 & 0.93 &  & 0.93 & 0.93 & 0.95 & 0.96 & 0.95 & 0.95 & 0.94 & 0.95 & 0.94 & 0.95 \\ 
  12 & 0.20 & 0.95 & 0.95 & 0.96 & 0.94 & 0.94 & 0.95 & 0.94 &  & 0.96 & 0.96 &  & 0.94 & 0.95 & 0.94 & 0.94 & 0.95 & 0.96 & 0.95 & 0.95 \\ 
  13 & 0.30 & 0.95 & 0.95 & 0.95 & 0.95 & 0.94 & 0.95 &  & 0.93 & 0.96 & 0.96 & 0.95 & 0.95 & 0.95 & 0.96 & 0.95 & 0.95 & 0.95 & 0.95 & 0.95 \\ 
  14 & 0.40 & 0.95 & 0.95 & 0.95 & 0.96 & 0.93 &  & 0.95 & 0.95 & 0.96 & 0.95 & 0.94 &  & 0.95 & 0.96 & 0.95 & 0.93 & 0.94 & 0.95 & 0.95 \\ 
  15 & 0.50 & 0.94 & 0.95 & 0.94 & 0.93 &  & 0.94 & 0.95 & 0.95 & 0.96 & 0.94 & 0.95 & 0.95 & 0.95 & 0.95 & 0.94 & 0.96 & 0.96 & 0.95 & 0.95 \\ 
  16 & 0.60 & 0.94 & 0.95 & 0.93 &  & 0.95 & 0.94 & 0.95 & 0.96 & 0.94 & 0.94 & 0.95 & 0.95 & 0.00 & 0.95 & 0.95 & 0.95 & 0.96 & 0.94 & 0.94 \\ 
  17 & 0.70 & 0.95 & 0.95 &  & 0.95 & 0.95 & 0.95 & 0.95 & 0.96 & 0.95 & 0.96 & 0.95 & 0.95 & 0.95 & 0.95 & 0.94 & 0.96 & 0.95 & 0.94 & 0.96 \\ 
  18 & 0.80 & 0.94 &  & 0.93 & 0.95 & 0.93 & 0.94 & 0.94 & 0.95 & 0.95 & 0.95 & 0.94 & 0.95 & 0.94 &  & 0.95 & 0.97 & 0.95 & 0.95 & 0.95 \\ 
  19 & 0.90 &  & 0.94 & 0.93 & 0.96 & 0.95 & 0.94 & 0.95 & 0.96 & 0.93 & 0.94 & 0.94 & 0.95 & 0.95 & 0.95 & 0.94 & 0.95 & 0.96 & 0.95 & 0.95 \\ 
   \hline
\end{tabular}
\end{table}

% latex table generated in R 4.4.2 by xtable 1.8-4 package
% Fri Jan  9 05:04:32 2026
\begin{table}[ht]
\centering
\begin{tabular}{rrrrrrrrrrrrrrrrrrrrr}
  \hline
 & ar & -0.9 & -0.8 & -0.7 & -0.6 & -0.5 & -0.4 & -0.3 & -0.2 & -0.1 & 0 & 0.1 & 0.2 & 0.3 & 0.4 & 0.5 & 0.6 & 0.7 & 0.8 & 0.9 \\ 
  \hline
1 & -0.90 & 0.94 & 0.95 & 0.96 & 0.95 & 0.95 & 0.95 & 0.95 & 0.95 & 0.96 & 0.95 & 0.96 & 0.96 & 0.95 & 0.95 & 0.96 & 0.95 & 0.94 & 0.95 &  \\ 
  2 & -0.80 & 0.95 & 0.95 & 0.93 & 0.96 & 0.95 &  & 0.95 & 0.95 & 0.95 & 0.95 & 0.94 & 0.95 & 0.94 & 0.95 & 0.95 & 0.97 & 0.95 &  & 0.93 \\ 
  3 & -0.70 & 0.95 & 0.95 & 0.95 & 0.95 & 0.96 & 0.95 & 0.95 & 0.96 & 0.94 & 0.94 & 0.96 & 0.95 & 0.96 & 0.94 & 0.95 & 0.94 &  & 0.92 & 0.92 \\ 
  4 & -0.60 & 0.94 & 0.94 & 0.94 & 0.96 & 0.95 & 0.95 & 0.00 & 0.95 & 0.93 & 0.94 & 0.95 & 0.95 & 0.95 & 0.95 & 0.94 &  & 0.93 & 0.95 & 0.95 \\ 
  5 & -0.50 & 0.95 & 0.96 & 0.95 & 0.95 & 0.95 & 0.94 & 0.96 & 0.95 & 0.94 & 0.95 & 0.95 & 0.95 & 0.94 & 0.93 &  & 0.91 & 0.95 & 0.94 & 0.96 \\ 
  6 & -0.40 & 0.96 & 0.95 & 0.95 & 0.93 & 0.96 & 0.94 & 0.95 &  & 0.95 & 0.94 & 0.95 & 0.95 & 0.93 &  & 0.93 & 0.94 & 0.95 & 0.95 & 0.94 \\ 
  7 & -0.30 & 0.94 & 0.95 & 0.95 & 0.94 & 0.95 & 0.93 & 0.95 & 0.95 & 0.95 & 0.94 & 0.95 & 0.93 &  & 0.93 & 0.94 & 0.95 & 0.95 & 0.96 & 0.95 \\ 
  8 & -0.20 & 0.95 & 0.95 & 0.95 & 0.95 & 0.95 & 0.95 & 0.94 & 0.95 &  & 0.95 & 0.92 &  & 0.94 & 0.95 & 0.93 & 0.94 & 0.94 & 0.94 & 0.94 \\ 
  9 & -0.10 & 0.96 & 0.94 & 0.95 & 0.95 & 0.95 & 0.94 & 0.94 & 0.95 & 0.95 & 0.93 &  & 0.92 & 0.94 & 0.95 & 0.95 & 0.96 & 0.96 & 0.96 & 0.95 \\ 
  10 & 0.00 & 0.94 & 0.95 & 0.94 & 0.95 & 0.95 & 0.94 & 0.94 & 0.95 & 0.91 &  & 0.93 & 0.94 & 0.93 & 0.96 & 0.95 & 0.95 & 0.95 & 0.95 & 0.95 \\ 
  11 & 0.10 & 0.95 & 0.95 & 0.94 & 0.95 & 0.95 & 0.95 & 0.95 & 0.91 &  & 0.92 & 0.95 & 0.95 & 0.95 & 0.95 & 0.94 & 0.94 & 0.94 & 0.96 & 0.96 \\ 
  12 & 0.20 & 0.95 & 0.95 & 0.94 & 0.95 & 0.95 & 0.94 & 0.93 &  & 0.91 & 0.93 &  & 0.94 & 0.94 & 0.97 & 0.95 & 0.95 & 0.94 & 0.96 & 0.94 \\ 
  13 & 0.30 & 0.95 & 0.95 & 0.94 & 0.95 & 0.93 & 0.95 &  & 0.92 & 0.95 & 0.95 & 0.94 & 0.95 & 0.94 & 0.95 & 0.95 & 0.95 & 0.95 & 0.95 & 0.95 \\ 
  14 & 0.40 & 0.94 & 0.94 & 0.95 & 0.93 & 0.93 &  & 0.93 & 0.95 & 0.93 & 0.96 & 0.95 &  & 0.95 & 0.95 & 0.95 & 0.94 & 0.95 & 0.94 & 0.95 \\ 
  15 & 0.50 & 0.95 & 0.96 & 0.94 & 0.93 &  & 0.94 & 0.95 & 0.96 & 0.96 & 0.94 & 0.95 & 0.95 & 0.95 & 0.95 & 0.95 & 0.95 & 0.96 & 0.96 & 0.94 \\ 
  16 & 0.60 & 0.95 & 0.95 & 0.93 &  & 0.94 & 0.95 & 0.94 & 0.96 & 0.94 & 0.95 & 0.95 & 0.95 & 0.00 & 0.95 & 0.95 & 0.94 & 0.95 & 0.95 & 0.95 \\ 
  17 & 0.70 & 0.94 & 0.90 &  & 0.94 & 0.94 & 0.95 & 0.94 & 0.95 & 0.96 & 0.95 & 0.95 & 0.95 & 0.95 & 0.95 & 0.95 & 0.97 & 0.95 & 0.96 & 0.95 \\ 
  18 & 0.80 & 0.94 &  & 0.94 & 0.95 & 0.96 & 0.95 & 0.95 & 0.95 & 0.95 & 0.95 & 0.94 & 0.96 & 0.95 &  & 0.96 & 0.95 & 0.94 & 0.95 & 0.95 \\ 
  19 & 0.90 &  & 0.95 & 0.95 & 0.95 & 0.96 & 0.95 & 0.95 & 0.95 & 0.95 & 0.96 & 0.94 & 0.95 & 0.94 & 0.95 & 0.96 & 0.94 & 0.95 & 0.96 & 0.95 \\ 
   \hline
\end{tabular}
\end{table}

% latex table generated in R 4.4.2 by xtable 1.8-4 package
% Fri Jan  9 03:51:44 2026
\begin{table}[ht]
\centering
\begin{tabular}{rrrrrrrrrrrrrrrrrrrrr}
  \hline
 & ar & -0.9 & -0.8 & -0.7 & -0.6 & -0.5 & -0.4 & -0.3 & -0.2 & -0.1 & 0 & 0.1 & 0.2 & 0.3 & 0.4 & 0.5 & 0.6 & 0.7 & 0.8 & 0.9 \\ 
  \hline
1 & -0.90 & 0.95 & 0.93 & 0.89 & 0.94 & 0.94 & 0.95 & 0.94 & 0.94 & 0.94 & 0.94 & 0.95 & 0.96 & 0.95 & 0.95 & 0.95 & 0.93 & 0.94 & 0.94 &  \\ 
  2 & -0.80 & 0.95 & 0.93 & 0.91 & 0.94 & 0.93 &  & 0.95 & 0.95 & 0.95 & 0.94 & 0.95 & 0.94 & 0.95 & 0.94 & 0.95 & 0.96 & 0.92 &  & 0.93 \\ 
  3 & -0.70 & 0.97 & 0.93 & 0.91 & 0.93 & 0.95 & 0.94 & 0.94 & 0.95 & 0.94 & 0.95 & 0.95 & 0.95 & 0.95 & 0.94 & 0.95 & 0.96 &  & 0.94 & 0.93 \\ 
  4 & -0.60 & 0.95 & 0.93 & 0.90 & 0.92 & 0.95 & 0.97 & 0.51 & 0.95 & 0.95 & 0.94 & 0.95 & 0.96 & 0.92 & 0.95 & 0.93 &  & 0.94 & 0.95 & 0.93 \\ 
  5 & -0.50 & 0.97 & 0.93 & 0.92 & 0.96 & 0.95 & 0.95 & 0.95 & 0.95 & 0.94 & 0.96 & 0.95 & 0.95 & 0.95 & 0.93 &  & 0.93 & 0.95 & 0.93 & 0.95 \\ 
  6 & -0.40 & 0.97 & 0.92 & 0.92 & 0.94 & 0.91 & 0.95 & 0.96 &  & 0.94 & 0.95 & 0.95 & 0.95 & 0.94 &  & 0.94 & 0.95 & 0.95 & 0.94 & 0.96 \\ 
  7 & -0.30 & 0.97 & 0.90 & 0.94 & 0.93 & 0.94 & 0.94 & 0.94 & 0.95 & 0.96 & 0.94 & 0.95 & 0.95 &  & 0.94 & 0.95 & 0.94 & 0.94 & 0.94 & 0.97 \\ 
  8 & -0.20 & 0.97 & 0.87 & 0.93 & 0.94 & 0.96 & 0.95 & 0.95 & 0.94 &  & 0.95 & 0.93 &  & 0.93 & 0.95 & 0.95 & 0.95 & 0.95 & 0.92 & 0.96 \\ 
  9 & -0.10 & 0.97 & 0.90 & 0.95 & 0.94 & 0.95 & 0.94 & 0.95 & 0.95 & 0.94 & 0.94 &  & 0.93 & 0.95 & 0.94 & 0.94 & 0.94 & 0.94 & 0.93 & 0.96 \\ 
  10 & 0.00 & 0.97 & 0.91 & 0.95 & 0.95 & 0.97 & 0.94 & 0.95 & 0.95 & 0.94 &  & 0.95 & 0.95 & 0.95 & 0.94 & 0.94 & 0.95 & 0.94 & 0.91 & 0.97 \\ 
  11 & 0.10 & 0.97 & 0.92 & 0.93 & 0.93 & 0.96 & 0.95 & 0.93 & 0.94 &  & 0.95 & 0.95 & 0.95 & 0.95 & 0.95 & 0.94 & 0.94 & 0.93 & 0.88 & 0.98 \\ 
  12 & 0.20 & 0.97 & 0.93 & 0.96 & 0.95 & 0.95 & 0.96 & 0.95 &  & 0.94 & 0.94 &  & 0.94 & 0.95 & 0.95 & 0.94 & 0.93 & 0.94 & 0.88 & 0.97 \\ 
  13 & 0.30 & 0.96 & 0.95 & 0.95 & 0.95 & 0.94 & 0.93 &  & 0.95 & 0.95 & 0.94 & 0.95 & 0.95 & 0.95 & 0.94 & 0.94 & 0.94 & 0.94 & 0.90 & 0.97 \\ 
  14 & 0.40 & 0.97 & 0.93 & 0.95 & 0.95 & 0.95 &  & 0.93 & 0.94 & 0.95 & 0.94 & 0.93 &  & 0.95 & 0.95 & 0.95 & 0.94 & 0.93 & 0.91 & 0.97 \\ 
  15 & 0.50 & 0.95 & 0.95 & 0.95 & 0.94 &  & 0.94 & 0.95 & 0.95 & 0.94 & 0.95 & 0.96 & 0.94 & 0.95 & 0.94 & 0.93 & 0.95 & 0.90 & 0.92 & 0.97 \\ 
  16 & 0.60 & 0.93 & 0.94 & 0.94 &  & 0.94 & 0.95 & 0.96 & 0.94 & 0.95 & 0.94 & 0.94 & 0.95 & 0.51 & 0.95 & 0.95 & 0.95 & 0.90 & 0.93 & 0.96 \\ 
  17 & 0.70 & 0.93 & 0.93 &  & 0.95 & 0.94 & 0.93 & 0.95 & 0.95 & 0.95 & 0.95 & 0.94 & 0.96 & 0.94 & 0.95 & 0.95 & 0.93 & 0.90 & 0.93 & 0.97 \\ 
  18 & 0.80 & 0.94 &  & 0.94 & 0.95 & 0.94 & 0.94 & 0.95 & 0.93 & 0.95 & 0.94 & 0.95 & 0.94 & 0.95 &  & 0.93 & 0.93 & 0.90 & 0.94 & 0.95 \\ 
  19 & 0.90 &  & 0.94 & 0.94 & 0.95 & 0.94 & 0.95 & 0.95 & 0.95 & 0.95 & 0.96 & 0.95 & 0.95 & 0.94 & 0.94 & 0.95 & 0.93 & 0.89 & 0.94 & 0.96 \\ 
   \hline
\end{tabular}
\end{table}

% latex table generated in R 4.4.2 by xtable 1.8-4 package
% Thu Jan  8 10:48:41 2026
\begin{table}[ht]
\centering
\begin{tabular}{rrrrrrrrrrrrrrrrrrrrr}
  \hline
 & ar & -0.9 & -0.8 & -0.7 & -0.6 & -0.5 & -0.4 & -0.3 & -0.2 & -0.1 & 0 & 0.1 & 0.2 & 0.3 & 0.4 & 0.5 & 0.6 & 0.7 & 0.8 & 0.9 \\ 
  \hline
1 & -0.90 & 0.90 & 0.93 & 0.98 & 0.97 & 0.96 & 0.94 & 0.93 & 0.94 & 0.97 & 0.97 & 0.96 & 0.96 & 0.93 & 0.91 & 0.95 & 0.94 & 0.93 & 0.87 &  \\ 
  2 & -0.80 & 0.99 & 0.94 & 0.95 & 0.95 & 0.92 &  & 0.92 & 0.94 & 0.97 & 0.96 & 0.97 & 0.95 & 0.94 & 0.95 & 0.97 & 0.96 & 0.81 &  & 0.79 \\ 
  3 & -0.70 & 0.97 & 0.92 & 0.97 & 0.96 & 0.90 & 0.94 & 0.93 & 0.93 & 0.88 & 0.95 & 0.91 & 0.94 & 0.97 & 0.96 & 0.93 & 0.83 &  & 0.81 & 0.89 \\ 
  4 & -0.60 & 0.93 & 0.95 & 0.93 & 0.96 & 0.95 & 0.95 & 0.00 & 0.93 & 0.94 & 0.92 & 0.94 & 0.95 & 0.97 & 0.88 & 0.82 &  & 0.74 & 0.87 & 0.95 \\ 
  5 & -0.50 & 0.95 & 0.99 & 0.94 & 0.95 & 0.95 & 0.92 & 0.98 & 0.93 & 0.96 & 0.97 & 0.96 & 0.90 & 0.91 & 0.81 &  & 0.80 & 0.82 & 0.96 & 0.92 \\ 
  6 & -0.40 & 0.94 & 0.96 & 0.92 & 0.96 & 0.94 & 0.94 & 0.97 &  & 0.92 & 0.94 & 0.90 & 0.94 & 0.78 &  & 0.77 & 0.80 & 0.88 & 0.90 & 0.90 \\ 
  7 & -0.30 & 0.90 & 0.92 & 0.90 & 0.94 & 0.96 & 0.96 & 0.93 & 0.91 & 0.94 & 0.92 & 0.87 & 0.77 &  & 0.77 & 0.89 & 0.91 & 0.89 & 0.96 & 0.98 \\ 
  8 & -0.20 & 0.95 & 0.96 & 0.93 & 0.95 & 0.92 & 0.91 & 0.92 & 0.96 &  & 0.89 & 0.71 &  & 0.76 & 0.86 & 0.86 & 0.85 & 0.93 & 0.97 & 0.96 \\ 
  9 & -0.10 & 0.94 & 0.94 & 0.99 & 0.97 & 0.91 & 0.94 & 0.92 & 0.88 & 0.86 & 0.81 &  & 0.69 & 0.90 & 0.94 & 0.95 & 0.93 & 0.92 & 0.93 & 0.93 \\ 
  10 & 0.00 & 0.94 & 0.94 & 0.95 & 0.91 & 0.92 & 0.96 & 0.90 & 0.85 & 0.79 &  & 0.81 & 0.89 & 0.96 & 0.95 & 0.92 & 0.92 & 0.91 & 0.97 & 0.93 \\ 
  11 & 0.10 & 0.95 & 0.90 & 0.95 & 0.90 & 0.93 & 0.92 & 0.81 & 0.75 &  & 0.76 & 0.85 & 0.94 & 0.96 & 0.95 & 0.91 & 0.95 & 0.90 & 0.97 & 0.93 \\ 
  12 & 0.20 & 0.91 & 0.94 & 0.91 & 0.93 & 0.91 & 0.90 & 0.72 &  & 0.84 & 0.92 &  & 0.96 & 0.97 & 0.93 & 0.97 & 0.97 & 0.94 & 0.97 & 0.93 \\ 
  13 & 0.30 & 0.92 & 0.93 & 0.91 & 0.93 & 0.87 & 0.83 &  & 0.75 & 0.84 & 0.94 & 0.95 & 0.94 & 0.98 & 0.92 & 0.95 & 0.96 & 0.94 & 0.94 & 0.96 \\ 
  14 & 0.40 & 0.94 & 0.92 & 0.95 & 0.90 & 0.75 &  & 0.79 & 0.86 & 0.94 & 0.94 & 0.97 &  & 0.94 & 0.95 & 0.98 & 0.93 & 0.91 & 0.91 & 0.93 \\ 
  15 & 0.50 & 0.86 & 0.89 & 0.84 & 0.78 &  & 0.84 & 0.87 & 0.93 & 0.97 & 0.97 & 0.95 & 0.95 & 0.90 & 0.93 & 0.93 & 0.94 & 0.91 & 0.95 & 0.95 \\ 
  16 & 0.60 & 0.95 & 0.88 & 0.84 &  & 0.80 & 0.88 & 0.90 & 0.96 & 0.96 & 0.96 & 0.98 & 0.94 & 0.00 & 0.98 & 0.97 & 0.93 & 0.95 & 0.97 & 0.96 \\ 
  17 & 0.70 & 0.91 & 0.76 &  & 0.74 & 0.92 & 0.93 & 0.95 & 0.97 & 0.94 & 0.95 & 0.96 & 0.95 & 0.98 & 0.94 & 0.95 & 0.92 & 0.94 & 0.97 & 0.97 \\ 
  18 & 0.80 & 0.81 &  & 0.86 & 0.96 & 0.98 & 0.93 & 0.94 & 0.93 & 0.87 & 0.96 & 0.93 & 0.95 & 0.95 &  & 0.91 & 0.94 & 0.96 & 0.94 & 0.89 \\ 
  19 & 0.90 &  & 0.94 & 0.93 & 0.96 & 0.96 & 0.94 & 0.98 & 0.95 & 0.97 & 0.95 & 0.95 & 0.97 & 0.93 & 0.95 & 0.92 & 0.96 & 0.96 & 0.96 & 0.96 \\ 
   \hline
\end{tabular}
\end{table}


To verify that such a method generates results consistent with theoretical confidence levels, we look at the results of simulating 1,000 realizations of an ARMA-Noise model with $T = 100,000$, $n = 2$, $\sigma_\eta^2 = 1$, $\sigma_\varepsilon^2 = 1$, each for combinations of $\phi$ and $\theta$ designated by the first two columns of the table below. For each realization, estimates were fit to simulated data and used to generate confidence regions at the $\alpha = .05$ level. The final column indicates the fraction of simulations that resulted in confidence regions that successfully captured the true time series parameters.


At first glance it appears that confidence region estimation under-performs for more extreme values of $\theta$. We look at simulations under such extreme conditions now for time series of length 1,000,000